\chapter{Referencia rápida}\label{ap:ref}

\section{Todos los tópicos}

{\small\begin{longtable}{@{}p{4.0cm} p{4.6cm} p{3.4cm} p{3.4cm}@{}}
\toprule
\textbf{Tópico} & \textbf{Tipo} & \textbf{Publica} & \textbf{Suscribe} \\
\midrule
\endhead
\topic{/cmd\_vel} & \cod{geometry\_msgs/Twist} & recorrido, teleop & S3 \\
\topic{/odom} & \cod{nav\_msgs/Odometry} & S3 (25~Hz) & recorrido \\
\topic{/imu/data\_raw} & \cod{sensor\_msgs/Imu} & S3 (50~Hz) & recorrido \\
\topic{/range/frente|izq|der} & \cod{sensor\_msgs/Range} & S3 (10~Hz) & recorrido \\
\topic{/bumpers} & \cod{std\_msgs/UInt8} & S3 (10~Hz) & recorrido \\
\topic{/battery} & \cod{sensor\_msgs/BatteryState} & S3 (1~Hz) & — \\
\topic{/base/estado} & \cod{std\_msgs/String} & S3 (2~Hz) & mision, recorrido \\
\topic{/base/comando} & \cod{std\_msgs/String} & recorrido, a mano & S3 \\
\topic{/base/respuesta} & \cod{std\_msgs/String} & S3 & — \\
\topic{/gps/fix} & \cod{sensor\_msgs/NavSatFix} & \cod{nmea\_navsat\_driver} (Pi, 1~Hz) & mision (hilera) \\
\topic{/brazo/joint\_states} & \cod{sensor\_msgs/JointState} & WROOM (20~Hz) & mision \\
\topic{/brazo/estado} & \cod{std\_msgs/String} & WROOM (2~Hz) & mision \\
\topic{/brazo/comando} & \cod{std\_msgs/String} & mision, a mano & WROOM \\
\topic{/brazo/respuesta} & \cod{std\_msgs/String} & WROOM & — \\
\topic{/brazo/joint\_cmd} & \cod{sensor\_msgs/JointState} & \cod{brazo\_ik} & WROOM, mision \\
\topic{/brazo/objetivo} & \cod{geometry\_msgs/PointStamped} & mision & \cod{brazo\_ik} \\
\topic{/brazo/ik\_estado} & \cod{std\_msgs/String} & \cod{brazo\_ik} & mision \\
\topic{/pinza/cierre} & \cod{std\_msgs/Float32} & mision & WROOM \\
\topic{/pinza/giro} & \cod{std\_msgs/Float32} & a mano & WROOM \\
\topic{/yolo/detecciones} & \cod{vision\_msgs/Detection2DArray} & \cod{detector\_node} & localizador, recorrido \\
\topic{/limon/posicion} & \cod{geometry\_msgs/PointStamped} & \cod{localizador\_node} & mision \\
\topic{/mision/estado} & \cod{std\_msgs/String} & mision & — \\
\topic{/recorrido/estado} & \cod{std\_msgs/String} & recorrido & — \\
\topic{/hilera/estado} & \cod{std\_msgs/String} & mision (1~Hz, con la hilera por GPS) & — \\
\bottomrule
\end{longtable}}

\section{Comandos de todos los días}

\begin{lstlisting}[language=consola]
# --- Conectarse y preparar
ssh r2d2@100.116.96.71
source /opt/ros/jazzy/setup.bash

# --- Después de prender la Pi (no arrancan solos)
bash ~/scripts/agente_base.sh     # agente de la base, en segundo plano
bash ~/scripts/vision.sh          # visión con YOLO v3

# --- Agentes micro-ROS a mano (una terminal cada uno)
source ~/microros_ws/install/local_setup.bash
ros2 run micro_ros_agent micro_ros_agent serial --dev /dev/serial/by-id/usb-Espressif_USB_JTAG_serial_debug_unit_*-if00
ros2 run micro_ros_agent micro_ros_agent serial --dev /dev/serial/by-id/usb-1a86_USB_Serial-if00-port0

# --- Base
ros2 topic pub --once /base/comando std_msgs/msg/String "{data: 'armar'}"
ros2 run teleop_twist_keyboard teleop_twist_keyboard --ros-args -p repeat_rate:=10.0
ros2 topic pub --once /base/comando std_msgs/msg/String "{data: 's'}"

# --- Brazo
ros2 topic pub --once /brazo/comando std_msgs/msg/String "{data: 'z'}"
ros2 topic pub --once /brazo/comando std_msgs/msg/String "{data: 'armar'}"
ros2 topic pub --once /brazo/joint_cmd sensor_msgs/msg/JointState "{name: [j1], position: [1.57]}"
ros2 topic pub --once /pinza/cierre std_msgs/msg/Float32 "{data: 50.0}"

# --- Visión a mano
MODELO=~/vision_estereo/modelos/limon_yolo26n_v3_ncnn_model CLASE=limon ~/vision_estereo/scripts/start_limon.sh

# --- Demo con el arbolito (capítulo de la demo; WROOM con el firmware 20, env brazo)
cd ~/scripts/27-aruco-arbol && python3 tools/recolectar.py

# --- Misión de ROS 2 (TF del 12/11)
ros2 launch mision mision.launch.py x:=0.1961 y:=-0.2694 z:=0.3833 roll:=0.0425 pitch:=0.2415 yaw:=1.5810 radio_limon:=0.025

# --- Firmware (en la notebook)
cd firmware/base && ~/.platformio/penv/bin/pio test -e native && ~/.platformio/penv/bin/pio run -e ros

# --- Grabar para la entrega
ros2 bag record /cmd_vel /odom /imu/data_raw /range/frente /bumpers /base/estado /limon/posicion /mision/estado /gps/fix /hilera/estado
\end{lstlisting}

\section{Estructura del repositorio}

\begin{lstlisting}
recolector_de_frutas/
|-- CLAUDE.md                 reglas del proyecto
|-- docs/
|   |-- INDICE.md             índice de documentos
|   |-- ARQUITECTURA.md       arquitectura del objetivo mínimo
|   |-- CONEXIONES.md         pines, potencia, autonomía, "antes de conectar"
|   |-- BITACORA.md           bitácora por fecha
|   |-- html/                 páginas interactivas (todo.html junta todas)
|   |-- ESTADO.md, PENDIENTES.md  estado de hoy y lo que falta
|   |-- manual/               este manual (manual.tex)
|   |-- fotos/                hojas de datos y fotos
|   |-- TPF R1/, TPF R2/      entregas
|-- firmware/
|   |-- lib/                  drivers y control compartidos
|   |-- base/                 ESP32-S3 (PlatformIO)
|   `-- brazo/                ESP32-WROOM (PlatformIO)
|-- ros2_ws/src/
|   |-- brazo_ik/             cinemática inversa
|   |-- mision/               misión + recorrido + TF cámara → brazo
|   `-- simulacion/           Gazebo Harmonic
`-- vision/entrenamiento_limon/   cuaderno de Colab del detector
\end{lstlisting}

\section{Documentos relacionados}
\begin{itemize}
  \item \archivo{docs/CONEXIONES.md}: la fuente de los pines, la potencia y la lista ``Antes de conectar''.
  \item \archivo{docs/html/todo.html}: todas las páginas de cableado en una, con cajas clickeables.
  \item \archivo{firmware/base/README.md}, \archivo{firmware/brazo/README.md}: detalle del firmware.
  \item \archivo{ros2\_ws/src/*/README.md}: detalle de cada paquete.
  \item \archivo{\textasciitilde/vision\_estereo/README.md}: detalle de la visión.
  \item \archivo{docs/BITACORA.md}: lo hecho día por día; \archivo{docs/DECISIONES.md}: por qué se decidió cada cosa.
  \item \archivo{docs/ESTADO.md} y \archivo{docs/PENDIENTES.md}: el estado más reciente (este manual es una foto al 19/11).
  \item Repositorio \cod{pruebas-esp32}: firmware 20 del brazo y scripts de la demo (\archivo{notas/ESTADO.md}).
\end{itemize}
